% year: תשע"ו
% type: מבחן
% semester: ב
% moed: ג
% source: exams/מבחנים/תשעו/מועד ג תשעו.pdf
% number: 2
% points: 10
% categories: true-false, continuity, func-limits

\begin{question}
\textbf{הוכיחו או הפריכו:} הפונקציה הבאה רציפה באפס:
\[ f(x)=\begin{cases} \frac{\sin(2x)}{\sqrt{x+4}-2} & x\neq 0\\ 5 & x=0\end{cases}. \]
\end{question}

\begin{hint}
הכפילו את המונה והמכנה בצמוד $\sqrt{x+4}+2$.
\end{hint}

\begin{solution}
\textbf{הטענה אינה נכונה.} $f$ רציפה ב-$0$ אם ורק אם $\lim_{x\to0}f(x)=f(0)=5$. נחשב את הגבול (עבור $x\ne0$, $x\ge-4$) בעזרת כפל בצמוד:
\[ \frac{\sin 2x}{\sqrt{x+4}-2}=\frac{\sin 2x\,(\sqrt{x+4}+2)}{(x+4)-4}=\frac{\sin2x}{x}\left(\sqrt{x+4}+2\right)=2\cdot\frac{\sin 2x}{2x}\cdot\left(\sqrt{x+4}+2\right). \]
לפי הגבול היסודי $\frac{\sin2x}{2x}\to1$, ובגלל רציפות השורש $\sqrt{x+4}+2\to4$. לפי אריתמטיקה של גבולות
\[ \lim_{x\to0}f(x)=2\cdot1\cdot4=8\neq5=f(0). \]
לכן $f$ \textbf{אינה רציפה} באפס (יש לה אי-רציפות סליקה; היא הייתה רציפה אילו הגדרנו $f(0)=8$).
\end{solution}
